% Generated by Claude Opus 5.5; no way I'm big brained enough to write these.
%
% In/out set diagram macros (\inOutDiagram, \inOutDiagramSides) for fig:InOutSets.
% Input from amspork3/amspork3_defs.tex ONLY.
% InOutSets layout: one tikz diagram per trace event constructor.
% \inOutDiagram[options]{constructor}{ObIn}{ObOut}{AwIn}{AwOut} (options: see "Manual overrides")
% The constructor box sits in the middle; Ob boxes above, Aw boxes below,
% "in" boxes on the left, "out" boxes on the right. Constructor and values are math.
% An empty value means \varnothing, drawn as a minimized "label = \varnothing" box.
% Every box hangs off its own corner of the constructor box by a 45-degree "leg"
% (same in both layouts, so no connection direction looks special).
% Breaks after "=" are manual only (see "Manual overrides"); by default every box is one line.
% Placement:
%   * \inOutDiagram: the "in" box's left edge and the "out" box's right edge sit just
%     outside the constructor's corners, so the legs land \vd@inset in from the boxes'
%     outer corners; a row too wide for that spreads outward symmetrically.
%   * \inOutDiagramSides: likewise vertically, per column.
\makeatletter
\newsavebox\vd@tmpbox
\newdimen\vd@W \newdimen\vd@H \newdimen\vd@gap \vd@gap=3mm
\newdimen\vd@leg \vd@leg=2mm     % horizontal = vertical extent of each corner leg
\newdimen\vd@inset \vd@inset=1.5mm % where the leg lands, measured from the box's outer corner
\newdimen\vd@a \newdimen\vd@b
\newdimen\vd@topShift \newdimen\vd@botShift
\newif\ifvd@breakTL \newif\ifvd@breakTR \newif\ifvd@breakBL \newif\ifvd@breakBR
\newcommand\vd@ifempty[3]{\if\relax\detokenize{#1}\relax #2\else #3\fi}
\tikzset{
  vertCtor/.style={draw=sporkMainColor, text=sporkMainColor, inner sep=3pt},
  vertVal/.style={draw=black!33, inner sep=2pt, align=left},
  vertEmpty/.style={draw=black!33, inner sep=1pt, align=left},
  vertStub/.style={draw=sporkMainColor, line width=1pt},
}
% Manual overrides, given as \inOutDiagram[options]{...} or \inOutDiagramSides[options]{...}:
%   tl break, tr break, bl break, br break = false (default) | true
%     (break that box after "="; \varnothing boxes never break)
%   top shift, bottom shift = <dimen> (\inOutDiagram only; replaces the automatic outward
%     spread of an overfull row; each box moves outward by this much)
\def\vd@strue{true}
\pgfkeys{/vd/.cd,
  tl break/.store in=\vd@optTL, tr break/.store in=\vd@optTR,
  bl break/.store in=\vd@optBL, br break/.store in=\vd@optBR,
  top shift/.store in=\vd@optTopShift, bottom shift/.store in=\vd@optBotShift}
% \vd@setbreak{option macro}{\vd@breakXXtrue}{\vd@breakXXfalse}
\newcommand\vd@setbreak[3]{\ifx#1\vd@strue #2\else #3\fi}
\newcommand\vd@setopts[1]{%
  \def\vd@optTL{false}\def\vd@optTR{false}\def\vd@optBL{false}\def\vd@optBR{false}%
  \let\vd@optTopShift\@empty \let\vd@optBotShift\@empty
  \pgfkeys{/vd/.cd,#1}%
  \vd@setbreak\vd@optTL\vd@breakTLtrue\vd@breakTLfalse
  \vd@setbreak\vd@optTR\vd@breakTRtrue\vd@breakTRfalse
  \vd@setbreak\vd@optBL\vd@breakBLtrue\vd@breakBLfalse
  \vd@setbreak\vd@optBR\vd@breakBRtrue\vd@breakBRfalse}
% \vd@ctorsize{constructor}: \vd@W, \vd@H = outer size of the constructor box
\newcommand\vd@ctorsize[1]{%
  \sbox\vd@tmpbox{\begin{tikzpicture}\node[vertCtor] {$\sporktr = #1$};\end{tikzpicture}}%
  \vd@W=\wd\vd@tmpbox \vd@H=\ht\vd@tmpbox \advance\vd@H by \dp\vd@tmpbox}
% \vd@measure{label}{value}{\vd@ifXX}: typesets the box as drawn into \vd@tmpbox
\newcommand\vd@measure[3]{%
  \sbox\vd@tmpbox{\begin{tikzpicture}\vd@box{}{tmp}{#1}{#2}#3\end{tikzpicture}}}
% \vd@spread{size A}{size B}{room}{shift dimen}: shift = max(0, A + B + gap - room) / 2
\newcommand\vd@spread[4]{%
  #4=#1 \advance#4 by #2 \advance#4 by \vd@gap \advance#4 by -#3
  \ifdim#4<0pt #4=0pt\fi #4=0.5#4}
% \vd@box{options}{node name}{label}{value}{break?}
\newcommand\vd@box[5]{%
  \vd@ifempty{#4}%
    {\node[vertEmpty, font=\scriptsize, #1] (#2) {$#3 = \varnothing$};}%
    {\node[vertVal, font=\small, #1] (#2) {#5{$#3 =$\\ $#4$}{$#3 = #4$}};}}
\newcommand\vd@ifTL{\ifvd@breakTL\expandafter\@firstoftwo\else\expandafter\@secondoftwo\fi}
\newcommand\vd@ifTR{\ifvd@breakTR\expandafter\@firstoftwo\else\expandafter\@secondoftwo\fi}
\newcommand\vd@ifBL{\ifvd@breakBL\expandafter\@firstoftwo\else\expandafter\@secondoftwo\fi}
\newcommand\vd@ifBR{\ifvd@breakBR\expandafter\@firstoftwo\else\expandafter\@secondoftwo\fi}
% The four 45-degree legs from the constructor box's corners.
\newcommand\vd@legs{%
  \draw[vertStub] (ctor.north west) -- ++(-\vd@leg,\vd@leg);
  \draw[vertStub] (ctor.north east) -- ++(\vd@leg,\vd@leg);
  \draw[vertStub] (ctor.south west) -- ++(-\vd@leg,-\vd@leg);
  \draw[vertStub] (ctor.south east) -- ++(\vd@leg,-\vd@leg);}
\newcommand{\inOutDiagram}[6][]{%
  \vd@setopts{#1}%
  \vd@ctorsize{#2}\advance\vd@W by 2\vd@leg \advance\vd@W by 2\vd@inset % room per row
  \vd@measure{\sporkObIn(\sporktr)}{#3}\vd@ifTL\vd@a=\wd\vd@tmpbox
  \vd@measure{\sporkObOut(\sporktr)}{#4}\vd@ifTR\vd@b=\wd\vd@tmpbox
  \vd@spread\vd@a\vd@b\vd@W\vd@topShift
  \vd@measure{\sporkAwIn(\sporktr)}{#5}\vd@ifBL\vd@a=\wd\vd@tmpbox
  \vd@measure{\sporkAwOut(\sporktr)}{#6}\vd@ifBR\vd@b=\wd\vd@tmpbox
  \vd@spread\vd@a\vd@b\vd@W\vd@botShift
  \ifx\vd@optTopShift\@empty\else \vd@topShift=\vd@optTopShift\relax\fi
  \ifx\vd@optBotShift\@empty\else \vd@botShift=\vd@optBotShift\relax\fi
  \begin{tikzpicture}[baseline=(ctor.base)]
    \node[vertCtor] (ctor) {$\sporktr = #2$};
    \vd@box{anchor=south west, at={($(ctor.north west)+(-\vd@leg-\vd@inset-\vd@topShift,\vd@leg)$)}}{tl}{\sporkObIn(\sporktr)}{#3}\vd@ifTL
    \vd@box{anchor=south east, at={($(ctor.north east)+(\vd@leg+\vd@inset+\vd@topShift,\vd@leg)$)}}{tr}{\sporkObOut(\sporktr)}{#4}\vd@ifTR
    \vd@box{anchor=north west, at={($(ctor.south west)+(-\vd@leg-\vd@inset-\vd@botShift,-\vd@leg)$)}}{bl}{\sporkAwIn(\sporktr)}{#5}\vd@ifBL
    \vd@box{anchor=north east, at={($(ctor.south east)+(\vd@leg+\vd@inset+\vd@botShift,-\vd@leg)$)}}{br}{\sporkAwOut(\sporktr)}{#6}\vd@ifBR
    \vd@legs
  \end{tikzpicture}%
}
% \inOutDiagramSides: same arguments, but the "in" quadrants sit to the left of the
% constructor and the "out" quadrants to the right (Ob upper, Aw lower).
% Trades horizontal space for vertical.
% \vd@topShift (left column) and \vd@botShift (right column) hold how far each box
% moves outward (vertically) if the column's two boxes would otherwise overlap.
\newcommand{\inOutDiagramSides}[6][]{%
  \vd@setopts{#1}%
  \vd@ctorsize{#2}\advance\vd@H by 2\vd@leg \advance\vd@H by 2\vd@inset % room per column
  \vd@measure{\sporkObIn(\sporktr)}{#3}\vd@ifTL\vd@a=\ht\vd@tmpbox \advance\vd@a by \dp\vd@tmpbox
  \vd@measure{\sporkAwIn(\sporktr)}{#5}\vd@ifBL\vd@b=\ht\vd@tmpbox \advance\vd@b by \dp\vd@tmpbox
  \vd@spread\vd@a\vd@b\vd@H\vd@topShift
  \vd@measure{\sporkObOut(\sporktr)}{#4}\vd@ifTR\vd@a=\ht\vd@tmpbox \advance\vd@a by \dp\vd@tmpbox
  \vd@measure{\sporkAwOut(\sporktr)}{#6}\vd@ifBR\vd@b=\ht\vd@tmpbox \advance\vd@b by \dp\vd@tmpbox
  \vd@spread\vd@a\vd@b\vd@H\vd@botShift
  \begin{tikzpicture}[baseline=(ctor.base)]
    \node[vertCtor] (ctor) {$\sporktr = #2$};
    \vd@box{anchor=north east, at={($(ctor.north west)+(-\vd@leg,\vd@leg+\vd@inset+\vd@topShift)$)}}{tl}{\sporkObIn(\sporktr)}{#3}\vd@ifTL
    \vd@box{anchor=north west, at={($(ctor.north east)+(\vd@leg,\vd@leg+\vd@inset+\vd@botShift)$)}}{tr}{\sporkObOut(\sporktr)}{#4}\vd@ifTR
    \vd@box{anchor=south east, at={($(ctor.south west)+(-\vd@leg,-\vd@leg-\vd@inset-\vd@topShift)$)}}{bl}{\sporkAwIn(\sporktr)}{#5}\vd@ifBL
    \vd@box{anchor=south west, at={($(ctor.south east)+(\vd@leg,-\vd@leg-\vd@inset-\vd@botShift)$)}}{br}{\sporkAwOut(\sporktr)}{#6}\vd@ifBR
    \vd@legs
  \end{tikzpicture}%
}
\makeatother
